% !TEX program = pdflatex
\documentclass[11pt,a4paper]{article}

% =========================================================================
% CONTROL DEL PROYECTO
% Proyecto: CasaScore (AlquilaPuno)
% Gestión de Proyectos TICS (EST315) -- Semestre 2026-II -- Team Sobrevivientes
% =========================================================================

% --- PAQUETES DE CONFIGURACIÓN ---
\usepackage[utf8]{inputenc}
\usepackage[T1]{fontenc}
\usepackage[spanish,es-tabla]{babel}
% Compatibilidad con babel-spanish reciente (flechas de TikZ y signo de porcentaje)
\AtBeginDocument{\shorthandoff{<>}\spanishplainpercent}
\usepackage{geometry}
\geometry{
    top=3.6cm,
    bottom=2.3cm,
    left=2.2cm,
    right=2.2cm,
    % PARTE MODIFICABLE: headheight debe ser mayor que la altura del logo más alto (2.56cm = 73pt).
    headheight=80pt,
    headsep=15pt,
    footskip=20pt
}

\usepackage{graphicx}
\usepackage{fancyhdr}
\usepackage[table]{xcolor}
\usepackage{titlesec}
\usepackage{booktabs}
\usepackage{tabularx}
\usepackage{array}
\usepackage{amsmath}
\usepackage{enumitem}
\usepackage{microtype}
\usepackage{amssymb}
% --- RECURSOS VISUALES: GRÁFICOS Y CÓDIGO DE COLORES ---
\usepackage{tikz}
\usetikzlibrary{positioning,calc}
\definecolor{GrisFila}{RGB}{246, 248, 251}
\definecolor{CuadCerca}{RGB}{150, 30, 30}
\definecolor{CuadSatisf}{RGB}{180, 130, 20}
\definecolor{CuadInfo}{RGB}{10, 45, 90}
\definecolor{CuadMonit}{RGB}{110, 115, 125}
\newcommand{\etiqueta}[2]{{\setlength{\fboxsep}{2pt}\colorbox{#1}{\textcolor{white}{\scriptsize\bfseries #2}}}}
\newcommand{\cerca}{\etiqueta{CuadCerca}{Gestionar de cerca}}
\newcommand{\satisf}{\etiqueta{CuadSatisf}{Mantener satisfecho}}
\newcommand{\informado}{\etiqueta{CuadInfo}{Mantener informado}}
\newcommand{\monit}{\etiqueta{CuadMonit}{Monitorear}}
\usepackage{needspace}
\usepackage{longtable}
\definecolor{Rama1}{RGB}{10, 45, 90}
\definecolor{Rama2}{RGB}{40, 95, 150}
\definecolor{Rama3}{RGB}{180, 130, 20}
\definecolor{Rama4}{RGB}{30, 110, 60}
\definecolor{Rama5}{RGB}{110, 60, 140}
\definecolor{RaciR}{RGB}{150, 30, 30}
\definecolor{RaciA}{RGB}{10, 45, 90}
\definecolor{RaciC}{RGB}{180, 130, 20}
\definecolor{RaciI}{RGB}{140, 145, 155}
\definecolor{NivAlto}{RGB}{30, 110, 60}
\definecolor{NivMedio}{RGB}{180, 130, 20}
\definecolor{NivBajo}{RGB}{150, 30, 30}
\usetikzlibrary{patterns}
\definecolor{SevBaja}{RGB}{30, 110, 60}
\definecolor{SevMedia}{RGB}{180, 130, 20}
\definecolor{SevAlta}{RGB}{150, 30, 30}
\usetikzlibrary{arrows.meta,shapes.geometric,fit,backgrounds}
\usepackage{listings}
\lstset{basicstyle=\ttfamily\scriptsize, frame=single, rulecolor=\color{AzulUNA!40}, backgroundcolor=\color{GrisFila}, columns=fullflexible, keepspaces=true, breaklines=true, xleftmargin=4pt, xrightmargin=4pt, literate={á}{{\'a}}1 {é}{{\'e}}1 {í}{{\'i}}1 {ó}{{\'o}}1 {ú}{{\'u}}1 {ñ}{{\~n}}1 {Á}{{\'A}}1}

\usepackage{hyperref}

% Logos: funciona a 2 o 3 niveles de profundidad dentro de docs/latex/
\graphicspath{{../../assets/img/}{../../../assets/img/}}

% --- PALETA INSTITUCIONAL FIEI / UNA PUNO ---
\definecolor{AzulUNA}{RGB}{10, 45, 90}
\definecolor{DoradoFIEI}{RGB}{180, 130, 20}
\definecolor{GrisOscuro}{RGB}{50, 50, 50}
\definecolor{AzulClaro}{RGB}{228, 236, 246}

% --- CONFIGURACIÓN DE LISTAS ---
\setlist{topsep=2pt, itemsep=2pt, parsep=1pt}

% --- COLUMNAS DE TABLA ---
\newcolumntype{L}{>{\raggedright\arraybackslash}X}
\newcolumntype{P}[1]{>{\raggedright\arraybackslash}p{#1}}
\newcolumntype{C}[1]{>{\centering\arraybackslash}p{#1}}
\newcolumntype{Q}[1]{>{\raggedleft\arraybackslash}p{#1}}
\newcolumntype{R}{>{\raggedleft\arraybackslash}X}
\renewcommand{\arraystretch}{1.25}

% --- FORMATO DE ENLACES ---
\hypersetup{
    colorlinks=true,
    linkcolor=AzulUNA,
    urlcolor=AzulUNA,
    citecolor=AzulUNA,
    pdftitle={Control del Proyecto -- CasaScore (AlquilaPuno)},
    pdfauthor={Team Sobrevivientes}
}

% --- FORMATO DE TÍTULOS ---
\titleformat{\section}{\color{AzulUNA}\normalfont\large\bfseries}{\thesection.}{0.7em}{}[\color{AzulUNA}\titlerule]
\titleformat{\subsection}{\color{AzulUNA}\normalfont\normalsize\bfseries}{\thesubsection.}{0.5em}{}
\titlespacing*{\section}{0pt}{9pt}{4pt}
\titlespacing*{\subsection}{0pt}{6pt}{2pt}

% --- ENCABEZADO Y PIE DE PÁGINA (fancyhdr) ---
\pagestyle{fancy}
\fancyhf{}

% ===== TAMAÑO DE LOS LOGOS (PARTE MODIFICABLE) =====
% logo_fiei.png  : height=1.45cm (lado izquierdo)
% logo_equipo.png: height=2.56cm (lado derecho; 25 % más grande porque la imagen tiene más borde blanco)
% Para cambiar el tamaño, modifique solo el valor height=... de cada \includegraphics.
% Si el logo más alto supera 2.6cm, aumente también headheight y top en la configuración de página.
\fancyhead[L]{%
    \includegraphics[height=1.45cm,keepaspectratio]{logo_fiei.png}%
}
\fancyhead[C]{%
    \scriptsize \textbf{\textcolor{AzulUNA}{UNIVERSIDAD NACIONAL DEL ALTIPLANO DE PUNO}}\\[2pt]
    \tiny \textcolor{GrisOscuro}{FACULTAD DE INGENIERÍA ESTADÍSTICA E INFORMÁTICA}\\[2pt]
    \scriptsize \textbf{\textcolor{AzulUNA}{Gestión de Proyectos TICS (EST315) \textbar\ Team Sobrevivientes}}\\[2pt]
    \footnotesize \textbf{\textcolor{DoradoFIEI}{CONTROL DEL PROYECTO}}%
}
\fancyhead[R]{%
    \includegraphics[height=2.56cm,keepaspectratio]{logo_equipo.png}%
}

\renewcommand{\headrulewidth}{0pt}
\renewcommand{\headrule}{}

\fancyfoot[L]{\scriptsize \textcolor{GrisOscuro}{CasaScore (AlquilaPuno)}}
\fancyfoot[C]{\small \thepage}
\fancyfoot[R]{\scriptsize \textcolor{GrisOscuro}{Puno, Perú -- Semestre 2026-II}}
\renewcommand{\footrulewidth}{0.4pt}

\begin{document}

\begin{center}
    {\Large \textbf{\textcolor{AzulUNA}{CONTROL DEL PROYECTO}}}\\[3pt]
    {\normalsize \textbf{CasaScore (AlquilaPuno)}}\\[4pt]
    \footnotesize \textbf{Fecha:} viernes 02/10/2026 (semana 6) \quad \textbar\quad \textbf{Versión:} 1.0 \quad \textbar\quad \textbf{Fase:} Monitoreo
\end{center}

\vspace{0.1cm}

\section{Datos Generales}
\noindent\begin{tabularx}{\textwidth}{@{}P{4.2cm} L@{}}
\toprule
\textbf{Proyecto:} & CasaScore (aplicación AlquilaPuno) \\
\textbf{Responsable:} & Ruth Karina Apaza Solis -- Analista de Planificación \\
\textbf{Proceso PM4DEV:} & Monitoreo y control (cronograma y presupuesto) \\
\textbf{Insumos:} & Línea base del cronograma, presupuesto, actas de sprint y registro de horas y gastos \\
\textbf{Fecha de corte:} & viernes 02/10/2026 (fin de la semana 6) \\
\bottomrule
\end{tabularx}

\section{Método}
El avance se mide con la gestión del valor ganado (EVM). Cada actividad de la EDT tiene un presupuesto (BAC); al corte se compara lo que debía estar hecho (PV), el valor de lo realmente hecho (EV $=$ \% real $\times$ BAC) y lo que costó (AC). El presupuesto del trabajo es S/ 11,108.07; la contingencia (S/ 555.40) se controla aparte. Los costos reales de este informe son simulados a partir del registro de horas y gastos del equipo.

\Needspace{14\baselineskip}
\section{Tablero de Control}
\begin{center}
\begin{tikzpicture}
\node[draw=SevAlta, fill=SevAlta!10, line width=1pt, rounded corners=4pt, minimum width=3.5cm, minimum height=2.0cm] at (0.00,0) {};\node[font=\scriptsize\bfseries, text=GrisOscuro] at (0.00,0.62) {CRONOGRAMA};\node[font=\LARGE\bfseries, text=SevAlta] at (0.00,0.05) {0.56};\node[font=\tiny, text=GrisOscuro, text width=3.3cm, align=center] at (0.00,-0.62) {SPI: atraso equivalente a 2 semanas};
\node[draw=SevMedia, fill=SevMedia!10, line width=1pt, rounded corners=4pt, minimum width=3.5cm, minimum height=2.0cm] at (3.80,0) {};\node[font=\scriptsize\bfseries, text=GrisOscuro] at (3.80,0.62) {COSTO};\node[font=\LARGE\bfseries, text=SevMedia] at (3.80,0.05) {0.96};\node[font=\tiny, text=GrisOscuro, text width=3.3cm, align=center] at (3.80,-0.62) {CPI: dentro del margen de 0.95};
\node[draw=SevAlta, fill=SevAlta!10, line width=1pt, rounded corners=4pt, minimum width=3.5cm, minimum height=2.0cm] at (7.60,0) {};\node[font=\scriptsize\bfseries, text=GrisOscuro] at (7.60,0.62) {AVANCE};\node[font=\LARGE\bfseries, text=SevAlta] at (7.60,0.05) {23.6\,\%};\node[font=\tiny, text=GrisOscuro, text width=3.3cm, align=center] at (7.60,-0.62) {del proyecto (plan: 42.2\,\%)};
\node[draw=Rama4, fill=Rama4!10, line width=1pt, rounded corners=4pt, minimum width=3.5cm, minimum height=2.0cm] at (11.40,0) {};\node[font=\scriptsize\bfseries, text=GrisOscuro] at (11.40,0.62) {EFECTIVO};\node[font=\LARGE\bfseries, text=Rama4] at (11.40,0.05) {646.92};\node[font=\tiny, text=GrisOscuro, text width=3.3cm, align=center] at (11.40,-0.62) {gastado de S/ 774.85 previstos al corte};
\end{tikzpicture}
\\[6pt]
\begin{tikzpicture}
\node[draw=SevMedia, fill=SevMedia!10, line width=1pt, rounded corners=4pt, minimum width=3.5cm, minimum height=2.0cm] at (0.00,0) {};\node[font=\scriptsize\bfseries, text=GrisOscuro] at (0.00,0.62) {RIESGOS ALTOS};\node[font=\LARGE\bfseries, text=SevMedia] at (0.00,0.05) {3};\node[font=\tiny, text=GrisOscuro, text width=3.3cm, align=center] at (0.00,-0.62) {R01, R02 y R07 en vigilancia};
\node[draw=AzulUNA, fill=AzulUNA!10, line width=1pt, rounded corners=4pt, minimum width=3.5cm, minimum height=2.0cm] at (3.80,0) {};\node[font=\scriptsize\bfseries, text=GrisOscuro] at (3.80,0.62) {CAMBIOS};\node[font=\LARGE\bfseries, text=AzulUNA] at (3.80,0.05) {1};\node[font=\tiny, text=GrisOscuro, text width=3.3cm, align=center] at (3.80,-0.62) {SC-01 rechazada};
\node[draw=SevAlta, fill=SevAlta!10, line width=1pt, rounded corners=4pt, minimum width=3.5cm, minimum height=2.0cm] at (7.60,0) {};\node[font=\scriptsize\bfseries, text=GrisOscuro] at (7.60,0.62) {CUARTOS};\node[font=\LARGE\bfseries, text=SevAlta] at (7.60,0.05) {4 / 15};\node[font=\tiny, text=GrisOscuro, text width=3.3cm, align=center] at (7.60,-0.62) {registrados para el piloto};
\node[draw=Rama4, fill=Rama4!10, line width=1pt, rounded corners=4pt, minimum width=3.5cm, minimum height=2.0cm] at (11.40,0) {};\node[font=\scriptsize\bfseries, text=GrisOscuro] at (11.40,0.62) {CONTINGENCIA};\node[font=\LARGE\bfseries, text=Rama4] at (11.40,0.05) {555.40};\node[font=\tiny, text=GrisOscuro, text width=3.3cm, align=center] at (11.40,-0.62) {sin uso al corte};
\end{tikzpicture}
\end{center}

\Needspace{30\baselineskip}
\section{Valor Ganado por Actividad}
Se muestran las actividades con valor planificado o ganado al corte. Valores en soles.

\vspace{0.1cm}
{\footnotesize
\setlength{\tabcolsep}{3pt}
\rowcolors{2}{white}{GrisFila}
\noindent\begin{tabularx}{\textwidth}{@{}C{0.7cm} L Q{1.4cm} Q{1.35cm} C{0.9cm} Q{1.35cm} Q{1.35cm} Q{1.4cm} Q{1.2cm}@{}}
\toprule
\rowcolor{AzulClaro}
\textbf{Cód.} & \textbf{Actividad} & \textbf{BAC} & \textbf{PV} & \textbf{\% real} & \textbf{EV} & \textbf{AC} & \textbf{SV} & \textbf{CV} \\
\midrule
1.1 & Gobernanza del equipo & 372.65 & 372.65 & 100\,\% & 372.65 & 372.65 & 0.00 & 0.00 \\
1.2 & Viabilidad y Acta de Constitución & 372.65 & 372.65 & 100\,\% & 372.65 & 380.10 & 0.00 & \textcolor{SevAlta}{-7.45} \\
1.3 & Planificación base & 496.86 & 496.86 & 40\,\% & 198.74 & 222.59 & \textcolor{SevAlta}{-298.12} & \textcolor{SevAlta}{-23.85} \\
1.4 & Planes de soporte & 372.65 & 372.65 & 10\,\% & 37.27 & 37.27 & \textcolor{SevAlta}{-335.38} & 0.00 \\
1.5 & Seguimiento y control & 1,061.03 & 83.04 & nivel & 83.04 & 83.04 & 0.00 & 0.00 \\
2.1 & Ficha técnica y requerimientos & 372.65 & 372.65 & 100\,\% & 372.65 & 372.65 & 0.00 & 0.00 \\
2.2 & Modelo de datos y arquitectura & 372.65 & 372.65 & 60\,\% & 223.59 & 234.77 & \textcolor{SevAlta}{-149.06} & \textcolor{SevAlta}{-11.18} \\
2.3 & Diseño de pantallas & 248.43 & 248.43 & 50\,\% & 124.22 & 129.18 & \textcolor{SevAlta}{-124.22} & \textcolor{SevAlta}{-4.97} \\
3.1 & Motor CASA SCORE & 496.86 & 496.86 & 100\,\% & 496.86 & 536.61 & 0.00 & \textcolor{SevAlta}{-39.75} \\
3.2 & Publicación de cuartos & 372.65 & 372.65 & 60\,\% & 223.59 & 237.01 & \textcolor{SevAlta}{-149.06} & \textcolor{SevAlta}{-13.42} \\
3.3 & Ranking y detalle (MVP Unidad I) & 1,034.45 & 443.34 & 0\,\% & 0.00 & 0.00 & \textcolor{SevAlta}{-443.34} & 0.00 \\
4.1 & Pruebas del motor & 248.43 & 212.94 & 0\,\% & 0.00 & 0.00 & \textcolor{SevAlta}{-212.94} & 0.00 \\
4.2 & Levantamiento de cuartos & 474.43 & 474.43 & 25\,\% & 118.61 & 130.47 & \textcolor{SevAlta}{-355.82} & \textcolor{SevAlta}{-11.86} \\
\midrule
\rowcolor{white}\multicolumn{3}{@{}l}{\textbf{Total al corte}} & \textbf{4,691.79} & & \textbf{2,623.86} & \textbf{2,736.33} & \textbf{\textcolor{SevAlta}{-2,067.93}} & \textbf{\textcolor{SevAlta}{-112.48}} \\
\bottomrule
\end{tabularx}
}

\Needspace{16\baselineskip}
\section{Indicadores}
{\small
\rowcolors{2}{white}{GrisFila}
\noindent\begin{tabularx}{\textwidth}{@{}C{1.1cm} P{3.6cm} L Q{2.0cm}@{}}
\toprule
\rowcolor{AzulClaro}
\textbf{Sigla} & \textbf{Indicador} & \textbf{Significado o fórmula} & \textbf{Valor} \\
\midrule
\textbf{PV} & Valor planificado & Costo del trabajo que debía estar hecho & \textbf{4,691.79} \\
\textbf{EV} & Valor ganado & Costo presupuestado del trabajo realmente hecho & \textbf{2,623.86} \\
\textbf{AC} & Costo real & Lo que realmente costó ese trabajo & \textbf{2,736.33} \\
\textbf{SV} & Variación del cronograma & EV $-$ PV & \textbf{\textcolor{SevAlta}{-2,067.93}} \\
\textbf{CV} & Variación del costo & EV $-$ AC & \textbf{\textcolor{SevMedia}{-112.48}} \\
\textbf{SPI} & Índice de cronograma & EV / PV & \textbf{\textcolor{SevAlta}{0.56}} \\
\textbf{CPI} & Índice de costo & EV / AC & \textbf{\textcolor{SevMedia}{0.96}} \\
\textbf{EAC} & Estimación al concluir & BAC / CPI & \textbf{\textcolor{SevMedia}{11,584.23}} \\
\textbf{ETC} & Estimación hasta concluir & EAC $-$ AC & \textbf{8,847.90} \\
\textbf{VAC} & Variación al concluir & BAC $-$ EAC & \textbf{\textcolor{SevMedia}{-476.16}} \\
\textbf{TCPI} & Desempeño requerido & (BAC $-$ EV) / (BAC $-$ AC) & \textbf{1.01} \\
\bottomrule
\end{tabularx}
}

\vspace{0.15cm}
\noindent\textbf{Lectura:} el equipo hizo el 56\,\% del trabajo previsto al corte; el valor ganado (S/ 2,623.86) equivale al valor planificado de la semana 4, por lo que el atraso es de \textbf{dos semanas}. El costo está controlado: por cada sol gastado se obtuvieron S/ 0.96 de trabajo. Si el CPI se mantiene, el proyecto costaría S/ 11,584.23, monto que aún entra en la línea base de costos (S/ 11,663.47) gracias a la contingencia.

\Needspace{20\baselineskip}
\section{Curva S al Corte}
\begin{center}
\begin{tikzpicture}[font=\scriptsize]
\draw[gray!20] (0,0.00) -- (13.00,0.00); \node[anchor=east, font=\tiny] at (0,0.00) {0};
\draw[gray!20] (0,0.92) -- (13.00,0.92); \node[anchor=east, font=\tiny] at (0,0.92) {2,000};
\draw[gray!20] (0,1.83) -- (13.00,1.83); \node[anchor=east, font=\tiny] at (0,1.83) {4,000};
\draw[gray!20] (0,2.75) -- (13.00,2.75); \node[anchor=east, font=\tiny] at (0,2.75) {6,000};
\draw[gray!20] (0,3.67) -- (13.00,3.67); \node[anchor=east, font=\tiny] at (0,3.67) {8,000};
\draw[gray!20] (0,4.58) -- (13.00,4.58); \node[anchor=east, font=\tiny] at (0,4.58) {10,000};
\draw[gray!20] (0,5.50) -- (13.00,5.50); \node[anchor=east, font=\tiny] at (0,5.50) {12,000};
\node[font=\tiny] at (0.00,-0.25) {1};
\node[font=\tiny] at (0.81,-0.25) {2};
\node[font=\tiny] at (1.62,-0.25) {3};
\node[font=\tiny] at (2.44,-0.25) {4};
\node[font=\tiny] at (3.25,-0.25) {5};
\node[font=\tiny] at (4.06,-0.25) {6};
\node[font=\tiny] at (4.88,-0.25) {7};
\node[font=\tiny] at (5.69,-0.25) {8};
\node[font=\tiny] at (6.50,-0.25) {9};
\node[font=\tiny] at (7.31,-0.25) {10};
\node[font=\tiny] at (8.12,-0.25) {11};
\node[font=\tiny] at (8.94,-0.25) {12};
\node[font=\tiny] at (9.75,-0.25) {13};
\node[font=\tiny] at (10.56,-0.25) {14};
\node[font=\tiny] at (11.38,-0.25) {15};
\node[font=\tiny] at (12.19,-0.25) {16};
\node[font=\tiny] at (13.00,-0.25) {17};
\draw[->] (0,0) -- (13.30,0) node[right, font=\tiny]{semana};
\draw[RaciR, dashed] (4.06,0) -- (4.06,5.5) node[above, font=\tiny, text=RaciR]{corte};
\draw[AzulUNA, very thick] (0,0) -- (0.81,0.256) -- (1.62,0.512) -- (2.44,1.109) -- (3.25,1.747) -- (4.06,2.150) -- (4.88,2.435) -- (5.69,2.778) -- (6.50,2.910) -- (7.31,3.190) -- (8.12,3.385) -- (8.94,3.772) -- (9.75,4.153) -- (10.56,4.368) -- (11.38,4.624) -- (12.19,4.985) -- (13.00,5.091);
\draw[Rama4, very thick] (0,0) -- (0.81,0.256) -- (1.62,0.512) -- (2.44,0.756) -- (3.25,0.985) -- (4.06,1.203);
\draw[SevAlta, very thick, dashed] (0,0) -- (0.81,0.256) -- (1.62,0.516) -- (2.44,0.781) -- (3.25,1.030) -- (4.06,1.254);
\draw[<->, GrisOscuro] (4.18,1.203) -- node[right, font=\tiny, text width=1.6cm]{SV = -2,067.93} (4.18,2.150);
\end{tikzpicture}
\\[2pt]
{\scriptsize \textcolor{AzulUNA}{\rule{12pt}{2pt}}~PV (línea base) \quad \textcolor{Rama4}{\rule{12pt}{2pt}}~EV (valor ganado) \quad \textcolor{SevAlta}{\rule{12pt}{2pt}}~AC (costo real)}
\end{center}

\Needspace{16\baselineskip}
\section{Solicitud de Cambio SC-01}
{\small
\noindent\begin{tabularx}{\textwidth}{@{}P{3.6cm} L@{}}
\toprule
\rowcolor{AzulClaro}
\multicolumn{2}{@{}l}{\textbf{Datos de la solicitud}} \\
\midrule
\textbf{Código y fecha:} & SC-01 \textperiodcentered\ presentada el martes 29/09/2026 (semana 6) \\
\textbf{Solicitante:} & Ruth Karina Apaza Solis -- Analista de Planificación \\
\textbf{Cambio propuesto:} & Aumentar la dedicación de cada integrante de 8 a 10 horas semanales desde la semana 7 hasta la semana 17, usando parte de las horas liberadas por la exclusión de Horus Sentinel. \\
\textbf{Justificación:} & Recuperar el atraso de dos semanas antes del piloto de la semana 13. \\
\bottomrule
\end{tabularx}
}

\vspace{0.2cm}
\noindent\textbf{Análisis de impacto}

\vspace{0.05cm}
{\small
\rowcolors{2}{white}{GrisFila}
\noindent\begin{tabularx}{\textwidth}{@{}P{2.4cm} L C{1.8cm}@{}}
\toprule
\rowcolor{AzulClaro}
\textbf{Área} & \textbf{Impacto} & \textbf{Efecto} \\
\midrule
Alcance & Sin cambio: se mantienen los requisitos del MVP. & \etiqueta{Rama4}{Neutro} \\
Cronograma & 88 horas adicionales del equipo (2 h $\times$ 4 integrantes $\times$ 11 semanas), equivalentes a 2.75 semanas de trabajo: recuperaría el atraso hacia la semana 13. & \etiqueta{Rama4}{Positivo} \\
Costo & S/ 1,061.50 más de personal valorizado (88 h $\times$ S/ 12.0625); el presupuesto total pasaría de S/ 12,829.82 a S/ 13,891.32. & \etiqueta{SevAlta}{Negativo} \\
Calidad & Más tiempo para pruebas y correcciones. & \etiqueta{Rama4}{Positivo} \\
Equipo y riesgos & Mayor carga junto con los exámenes parciales (semana 9) y las demás asignaturas; aumenta el riesgo de ausencias y de baja calidad por cansancio. & \etiqueta{SevAlta}{Negativo} \\
\bottomrule
\end{tabularx}
}

\vspace{0.2cm}
\noindent\textbf{Decisión del equipo (sponsor):} \etiqueta{SevAlta}{RECHAZADA} por unanimidad en la reunión de control del martes 29/09/2026. El aumento de costo y la carga académica pesan más que la recuperación del plazo. Se mantiene la dedicación de 8 horas semanales.

\vspace{0.1cm}
\noindent\textbf{Acción alternativa aprobada:} plan de recuperación sin horas adicionales: ejecutar en paralelo el ranking y la revisión de los documentos de planificación, levantar cuartos los sábados dentro de las 8 horas, y concentrar la revisión de documentos en las semanas 9 y 10, cuando el desarrollo se detiene por exámenes.

\Needspace{10\baselineskip}
\section{Registro de Cambios}
{\small
\noindent\begin{tabularx}{\textwidth}{@{}C{1.1cm} C{1.6cm} L P{2.0cm} C{1.6cm} C{1.6cm}@{}}
\toprule
\rowcolor{AzulClaro}
\textbf{Código} & \textbf{Fecha} & \textbf{Cambio} & \textbf{Solicitante} & \textbf{Estado} & \textbf{Decisión} \\
\midrule
SC-01 & 29/09/2026 & Aumentar la dedicación de 8 a 10 horas semanales & Ruth K. Apaza Solis & \etiqueta{SevAlta}{Rechazada} & 29/09/2026 \\
\bottomrule
\end{tabularx}
}

\Needspace{16\baselineskip}
\section{Informe de Variación del Presupuesto}
Comparación entre el costo previsto al corte según la línea base y el costo real registrado.

\vspace{0.1cm}
{\footnotesize
\setlength{\tabcolsep}{3.5pt}
\rowcolors{2}{white}{GrisFila}
\noindent\begin{tabularx}{\textwidth}{@{}P{2.7cm} Q{1.6cm} Q{1.6cm} Q{1.5cm} Q{1.1cm} L@{}}
\toprule
\rowcolor{AzulClaro}
\textbf{Rubro} & \textbf{Previsto al corte} & \textbf{Real al corte} & \textbf{Variación} & \textbf{\%} & \textbf{Causa} \\
\midrule
Personal regular & 3,222.28 & 1,609.81 & -1,612.47 & -50\,\% & Menos horas registradas que las planificadas. \\
Horas extra & 215.06 & 0.00 & -215.06 & -100\,\% & Aún no se realizó la jornada nocturna prevista. \\
Equipos & 479.60 & 479.60 & 0.00 & 0\,\% & Depreciación según el tiempo transcurrido. \\
Servicios & 282.02 & 282.02 & 0.00 & 0\,\% & Gasto fijo mensual. \\
Productos de salud & 163.83 & 120.50 & -43.33 & -26\,\% & Compras menores a las previstas. \\
Alimentación & 245.86 & 198.40 & -47.46 & -19\,\% & Menos reuniones presenciales. \\
Transporte & 83.14 & 46.00 & -37.14 & -45\,\% & Menos visitas de campo por el atraso del levantamiento. \\
\midrule
\rowcolor{white}\textbf{Total} & \textbf{4,691.79} & \textbf{2,736.33} & \textbf{-1,955.46} & \textbf{-42\,\%} & \\
\bottomrule
\end{tabularx}
}

\vspace{0.15cm}
\noindent El gasto real es menor al previsto porque se hizo menos trabajo, no porque se ahorrara: por eso la variación del presupuesto debe leerse junto con el valor ganado. El gasto en efectivo al corte es de S/ 646.92 frente a S/ 774.85 previstos.

\Needspace{10\baselineskip}
\section{Acciones Correctivas}
{\small
\rowcolors{2}{white}{GrisFila}
\noindent\begin{tabularx}{\textwidth}{@{}L P{3.4cm} C{1.6cm}@{}}
\toprule
\rowcolor{AzulClaro}
\textbf{Acción} & \textbf{Responsable} & \textbf{Hasta} \\
\midrule
Construir el ranking y el detalle en paralelo con la revisión de documentos & Eliseo Tarqui Ajahuana & 09/10 \\
Completar el levantamiento de cuartos con visitas los sábados & Clyde N. Paricahua Pari & Semana 12 \\
Revisar y firmar los documentos de planificación en las semanas de exámenes & Viviana M. Quisbert Quispe & Semana 10 \\
Registrar las horas reales cada día en el Issue de cada tarea & Todos & Permanente \\
\bottomrule
\end{tabularx}
}

\Needspace{24\baselineskip}
\section{Firmas}
\noindent En señal de conformidad con el presente documento, los integrantes del equipo firman a continuación:

\vspace{1.5cm}
\begin{center}
\begin{tabular}{c c c}
    \parbox{6.5cm}{\centering \rule{6.0cm}{0.6pt}\\[5pt] \textbf{Clyde Neil Paricahua Pari}\\[2pt] {\footnotesize Project Manager}}
    & \hspace{0.8cm} &
    \parbox{6.5cm}{\centering \rule{6.0cm}{0.6pt}\\[5pt] \textbf{Ruth Karina Apaza Solis}\\[2pt] {\footnotesize Analista de Planificación}}
    \\[2.0cm]
    \parbox{6.5cm}{\centering \rule{6.0cm}{0.6pt}\\[5pt] \textbf{Eliseo Tarqui Ajahuana}\\[2pt] {\footnotesize Especialista Técnico / Arquitectura}}
    & \hspace{0.8cm} &
    \parbox{6.5cm}{\centering \rule{6.0cm}{0.6pt}\\[5pt] \textbf{Viviana Mary Quisbert Quispe}\\[2pt] {\footnotesize Gestora de Comunicación, Riesgos y Documentación}}
\end{tabular}
\end{center}
\end{document}
